\section{A general upper bound for $C_k$}\label{sec:Ck}

% \commentt{Change section title? Also order of sections?}

\commentt{I started writing comments in several places. But maybe it is best that I summarize here what are the changes I am planning to make (only in the way things are presented, not in the actual technical details.}

\commentt{
In terms of the order of things, after the definition of $\ell(F)$ (and a short discussion of it), the first algorithm for $F$ (with sampled vertices only) will be presented and analyzed (including the main claim regarding hitting tuples). The corollary for $C_k$ will appear afterwards. Then we shall give the odd-$k$  algorithm for $C_k$, with its analysis (including the edge-sampling).}

\commentt{
I think that the argument regarding the existence of sufficiently many light-vertex-disjoint copies is similar to things we did before, so it can appear in the proofs of the Theorems. 
}

In this we section describe our algorithms for testing 
$C_k$-freeness, for even and odd $k$, and more generally for $F$-freeness. The proofs of correctness follow.
\commentt{maybe use $k$ for $|V(F)|$? Also, do we want to say something here about the running times?}

\newcommand{\VC}{\mathcal{VC}}
\commentt{The problem with the notation $\calC$ is that we used it a lot for a set of edge-disjoint cycles. So maybe $\VC$. Also, using $C$ might not be great, since it might be confused with $C_k$}
\begin{definition}\label{def:ell}
    For a graph $F = (V_F, E_F)$ let $\mathcal{C}(F)$ denote the set of vertex covers of $F$. For a vertex cover $C$ of $F$ we denote by $\mathcal{C'}(C)$ the set of vertex covers of $F$ which are subsets of $C$.
    We define $\ell(F) = max_{C \in \mathcal{C}(F)} \left\{min_{B \in \mathcal{C'}(C)} \left(|B|\right)\right\}$.
\end{definition}
\commentt{Maybe add an explanation about this notion, in particular that it may be much larger than a minimum vertex cover, e.g., for a star?}

\begin{claim}
Let $F = (V_F, E_F)$ and $G = (V_G, E_G)$ be a pair of graphs such that $F$ is a subgraph of $G$. 
Let $C$ be a vertex cover of $G$.
There exists a set of $\ell = \ell(F)$ vertices $v_1, \ldots, v_\ell$ in $C$ such that $E_F \subseteq \bigcup_{i\in [\ell]} \Gamma(v_i)$ where $\Gamma(v)$ denotes the set of edges incident to $v$.  
\commentt{maybe use $E(v)$ because $\Gamma(v)$ is usually used for the neighbors of $v$. Also, I didn't see where we refer to this claim, so I guess it is implicit and maybe we don't need to state it?}
\end{claim}

\begin{proof}
Since $C$ covers the edges of $G$ is also covers the edges of its subgraphs and in particular the edges of $F$. Therefore $D = C\cap V_F$ is a vertex cover of $F$. Consider $B\in \mathcal{C'}(D)$ of minimum size. Such $B$ is a vertex cover of $F$ and its size is at most $\ell(F)$.
Since $B$ is a vertex cover of $F$ it follows that $E_F \subseteq \cup_{v\in B} \Gamma(v)$, as desired.  
\end{proof}


\commentt{shouldn't the algorithms have different names? Or actually, are we keeping only the first one? And if so, are some of the claims later on redundant?}

\commentt{Also, to be consistent with previous parts, I suggest to replace $t$ with $\theta_0$. Also, there is the notion here of light vertices which is only explained later }

\commentt{I wonder how much effort we should put into trying to make things consist with the previous sections}

\commentt{The issue of getting $m$ as input should be discussed somewhere.}

\medskip
    \textsf{Test-$F$-freeness for general $F$}($n, m, \alpha, \epsilon$)  \commentt{Maybe remove ``for general $F$''}
\begin{enumerate}
%\item Obtain $\alpha^*$ (see Claim~\ref{lem:alpha}) and set $t = 48\alpha^*/\epsilon$.
\item Set $t = 4\alpha/\epsilon$.
\item Sample u.a.r. $s = $$\Theta\left(|V(F)|^{2+1/\ell(F)} \cdot m \cdot \left(\frac{\alpha}{m }\right)^{1/\ell(F)} \cdot \left(\frac{1}{\epsilon}\right)^{1+2/\ell(F)}\right)$ vertices (see Definition~\ref{def:ell}).
\item Reveal the neighborhood of all the light vertices (according to $t$) in the sample. 
\item Reject iff there is a witness for $F$ in the resulting subgraph
\end{enumerate}

\medskip
    \textsf{Test-$F$-freeness for general $F$}($n, m, \alpha, \epsilon$)
\begin{enumerate}
%\item Obtain $\alpha^*$ (see Claim~\ref{lem:alpha}) and set $t = 48\alpha^*/\epsilon$.
\item Set $t = 4\alpha/\epsilon$.
\item Sample u.a.r. $h = \Theta\left(|E(F)|^{1+1/|E(F)|} \cdot m^{1-1/|E(F)|} \cdot \left(\frac{1}{\epsilon}\right)^{1/|E(F)|}\right)$ edges from $H(G, t)$ (see Claim~\ref{clm:sample}).
\item Reject iff there is a witness for $F$ in the resulting subgraph
\end{enumerate}


\textsf{Test-$C_k$-freeness when $k$ is odd}($n, m, \alpha, \epsilon$)
\begin{enumerate}
% \item Obtain $\alpha^*$ (see Claim~\ref{lem:alpha}) and set $t = 48\alpha^*/\epsilon$.
\item Set $t = 4\alpha/\epsilon$.
\item Sample u.a.r. $s_1 = \Theta \left(k \cdot m^{1-\frac{2}{k-1}} \cdot (1/\alpha)^{1-\frac{4}{k-1}} \cdot \epsilon^{1-\frac{6}{k-1}}\right)$ edges from the graph. For each edge in the sample,  reveal the neighborhood of both endpoints if both are light. \commentt{I guess this should be the same as in the previous algorithm}
\item If $k > 3$, then sample u.a.r. $s_2 = O\left(k \cdot n \cdot (\alpha^2/m)^{\frac{2}{k-1}} \cdot \left( \frac{1}{\epsilon}\right)^{\frac{6}{k-1}} \right)$ vertices and reveal the neighborhood of all the light vertices in the sample. 
\item Reject iff there is a witness for $F$ in the resulting subgraph.
\end{enumerate}

Before we prove the correctness of the algorithms we shall state several useful definitions and lemmas. 

\commentt{Some of the following is common to all our algorithms. Ideally, I would suggest to  move a general version to the preliminaries and then use it everywhere, or maybe it will be simpler to just say here, the similarly to previous sections, there exists a set $\mathcal{F}$ of at least $\eps m/(2k)$ edge-disjoint copies of $F$ in $G$ (where $k = |V(F)|$ such that no copy includes an edge between two heavy vertices?}

\begin{definition}(\cite{})
For a graph $G = (V, E)$ and a threshold $t$ we define the set of {\em heavy vertices with respect to $t$} as $H_t(G) = \{v\in V: d(v) > t\}$ 
and the set of {\em light vertices with respect to $t$} as $L_t(G) = V \setminus H_t(G)$.
\end{definition}
When $G$ is clear from the context we may simply use $H_t$ and $L_t$. 
Using the definition of $H_t(G)$, we next define the graph $H(G, t)$ which is defined w.r.t. $G$ and a threshold $t$.

\begin{definition}[\bf The undirected graph $H(G, t)$](\cite{})
For a graph $G = (V, E)$ and a threshold $t$, the graph $H(G, t)$ is an undirected graph defined as follows.
The set of vertices of $H(G, t)$ is $V$ and the set of edges of $H(G, t)$ is $E(G)\setminus \{\{u, v\} : u\in H_t(G) \text{ and } v \in H_t(G)\}$.
Namely, $H(G, t)$ is the graph $G$ after removing the edges for which both endpoints are heavy with respect to $t$. 
\end{definition}

\begin{claim}\label{clm:arb}(\cite{})
For a graph $G$, $\alpha' \geq \alpha(G)$ and $\eta\in (0,1]$ it holds that $|E(H(G, \alpha'/\eta))| \geq (1-2\eta)|E(G)|$.
\end{claim}


\begin{claim}\label{lem:alpha}(\cite{})
There exists an algorithm that outputs $\alpha^*$ for which the following holds with probability at least $5/6$:
\begin{enumerate}
\item $|E(H(G, t))| \geq (1-(\eps/6))m$ where $t = 48\alpha^*/\epsilon$,
\item $\alpha^* \leq 2 \alpha(G)$.
\end{enumerate}
\end{claim} 

We say that a pair of subgraphs are {\em light-vertex-disjoint} if they do not share any light vertex (where a light vertex is defined w.r.t. some threshold, depending on the context). 

\begin{claim}\label{clm:disjoint}
Let $G$ be a graph that is $\epsilon$-far from being $F$-free. Let $t$ be a threshold such that $|E(H(G, t))| \geq (1-(\epsilon/2)) |E(G)|$. There exist at least $\epsilon m/ (2t |V(F)|)$ light-vertex-disjoint copies of $F$ in $H(G, t)$ (where a light vertex is defined w.r.t. $t$).
\end{claim}
\begin{proof}
Consider the following process that removes edges from $|E(H(G, t))|$. At each step, if there is a copy of $F$ in the graph, it removes all the edges adjacent to any one of its light vertices. The process terminates after all such edges are removed. Clearly, when the process terminates it has to remove at least $\epsilon m /2$ edges (otherwise we reach a contradiction to the fact that $G$ is $\epsilon$-far from being $F$-free).
At each step, the process removes at most $t \cdot |F|$ edges. Thus, we obtain the desired result. 
\end{proof}

\commentt{What appears below is almost the same as a procedure we already have for sampling edges. But there are some differences. There is also the issue of expected vs maximum}

\medskip
We slightly modify an algorithm from~\cite{} to obtain Claim~\ref{clm:sample}. 

\textsf{Sample a uniform edge from $H(G, t)$}
\begin{enumerate}
\item While do:
\begin{enumerate}
\item Sample u.a.r. a vertex from $V(G)$.
\item Pick $j\in [t]$ and toss a fair coin $b\in \{\rm{Heads}, \rm{Tails}\}$.
\item If $v\in L_t(v)$ and $v$ has a $j$-th neighbour, $u$ then:
\begin{enumerate}
\item If $u \in H_t(G)$ then return $\{u,v\}$. 
\item If $u \in L_t(G)$ and $b=\rm{Heads}$ then return $\{u,v\}$.
\end{enumerate}
\end{enumerate}
\end{enumerate}


\begin{claim}\label{clm:sample}
There exists an algorithm that samples a uniform edge from $H(G, t)$.
If $t$ is such that $|E(H(G, t))| \geq |E(G)|/2$ then the expected running time of the algorithm is $O(t/\bar{d}(G))$. 
\end{claim}



\begin{claim}\label{clm:sampleedges}
Let $S$ be a set of $\beta m$ edges in $H(G, t)$, where $\beta \in (0, 1]$. If we sample $\log n/\beta  \leq y \leq x/2$ edges uniformly from $H(G, t)$, then with high probability we sample at least $c \beta y$ different edges from $x$, where $c$ is an absolute constant. 
\end{claim}

\begin{proof}
    Let $X_i$ denote the indicator variable for the event that the $i$-th sample is an edge from $S$. Since the edges are sampled uniformly from $H(G, t)$, $\Ex(X_i) = \beta$.
    Thus, by multiplicative Chernoff's bound, with high probability $\sum_i X_i = b \beta y$, where $b$ is a constant. 
    Let $Y_i$ be the indicator variable for the event that the $i$-th sample is a new edge from $S$ (with respect the the $(i-1)$ first samples).
    We consider, without loss of generality, the sampling process as if it first decides the values of $X_i$ for each sample $i$ and then sample the edges one by one accordingly.
    Consider the $j$-th sample.
    Given the values of $X_i$ for every $i$ and the values of the samples for each $i < j$, the probability that the $j$-th sample is a new edge, conditioned on the event that $X_j = 1$ is constant (since $y \leq x/2)$. Therefore, with high probability we sample $c \beta y$ different edges from $S$ for some absolute constant $c$. 
\end{proof}

Let $X$ be a set of elements of size $x$ and let $C=\{(v^j_1, \ldots, v^j_\ell)\}_{j\in [|C|]}$ be a set of $\ell$-tuples of elements of $X$ such that each $v\in X$ appears at most once in at most one of the $\ell$-tuples of $C$ (namely, the $\ell$-tuples are disjoint).

Let $S = \{u_1, \ldots u_s\}$ be a set of $s$ elements chosen independently and u.a.r. from $X$. We say that an $\ell$-tuple $(i_1, \ldots i_\ell)$ where $i_j \in s$ for each $j\in [\ell]$ is good w.r.t. $C$ if $(u_{i_1}, \ldots u_{i_\ell})$ belongs to $C$.

We prove the following lemma.

\begin{claim}\label{clm:ver-num}
If $s \geq 16 \ell x / (|C|^{1/\ell})$ then with high constant probability there exists an $\ell$-tuple of $S$ which is good w.r.t. $C$.
\end{claim}

\begin{proof}
Let $\theta$ denote the number of different $\ell$-tuples of $S$.
\commentt{Better not to use $\theta$ for this since we use $\theta_0$ and $\theta_1$ a lot for thresholds. And maybe $\mathcal{L}$ or $\mathcal{T}$ for the tupes?}
Consider all the $\ell$-tuples of $S$ ordered according to some arbitrary order. 
For every $i\in [\theta]$, define $\chi_i$ to be the indicator variable for the event that the $i$-th tuple is good w.r.t. $C$.

\commentt{can we use one of the claims we already have (which we moved from the preliminaries to the appendix) or it doesn't fit?}
We next upper bound the probability that
$\sum_{i \in [\theta]} \chi_{i}=0$.
Let $\bar{\chi}_{i} = \chi_{i} - \mu$, so that $\Ex[\bar{\chi}_{i}] = 0$.
By Chebyshev's inequality,
\begin{equation}\label{eq:cheb-Ck}
\Pr\left[\sum_{i\in [\theta]}\chi_{i}=0\right] \leq \frac{\Var\left[\sum_{i\in [\theta]}\chi_{i}\right]}{\left(\Ex\left[\sum_{i\in [\theta]}\chi_{i}\right]\right)^2}
= \frac{\sum_{i,j} \Ex[\bar{\chi}_{i}\cdot \bar{\chi}_{j}]}{\theta^2\cdot \mu^2} \;.
\end{equation}

We first calculate 
\begin{equation}
\Ex\left[\sum_{i\in [\theta]}\chi_{i}\right] \geq \theta \cdot \mu \geq {s \choose \ell } \cdot \ell! \cdot |C| \cdot \left(\frac{1}{x}\right)^{\ell} 
= \frac{s!}{(s-\ell)!} \cdot |C| \cdot \left(\frac{1}{x}\right)^{\ell} \geq (s/2)^\ell \cdot |C| \cdot \left(\frac{1}{x}\right)^{\ell}  \geq (8\ell)^\ell 
\end{equation}

We next break the sum $\sum_{i,j} \Ex[\bar{\chi}_{i}\cdot \bar{\chi}_{j}]$ into several sub-sums as follows.
We say that $(i,j) \in [\theta]\times [\theta]$ is in $S_t$ if the number of elements (namely, indexes of samples) that belong both to the $i$-th tuple and the $j$-th tuple is exactly $t$ (namely the intersection of the tuples is of size $t$ when they are viewed as sets).  

If $(i,j) \in S_0$ then $\chi_{i}$ and $\chi_{j}$ are independent and so
\begin{equation}\label{eq:cheb-case1}
\Ex[\bar{\chi}_{i}\cdot \bar{\chi}_{j}] =  \Ex[\bar{\chi}_{i}]\cdot \Ex[\bar{\chi}_{j}] = 0\;.
\end{equation}
% where the last iequality follows from the definition of $\bar{\chi}_{i,j}$.
Hence, the pairs in $S_0$ do not contribute anything to the right-hand side of Equation~\eqref{eq:cheb-Ck}.


If $i = j$ then 
\begin{equation}\label{eq:cheb-Ck-case2}
\Ex[\bar{\chi}_{i}\cdot \bar{\chi}_{j}] = \Ex[(\bar{\chi}_{i})^2]
\leq \Ex[(\chi_{i})^2] = \Ex[\chi_{i}] = \mu \;,
\end{equation}
where the second inequality follows from the definition of $\bar{\chi}_{i}$.

Given that $(i,j) \in S_t$ it holds that  

\begin{equation}\label{eq:cheb-Ck-case2}
\Ex[\bar{\chi}_{i}\cdot \bar{\chi}_{j}] 
\leq \Ex[\chi_{i} \chi_{j}] \leq  |C| \cdot \left(\frac{1}{x}\right)^{2\ell -t} \;.
\end{equation}
Since there are at most ${\ell \choose t} \cdot s^{2\ell -t}$ pairs in $S_t$ we obtain that 
\begin{equation}\label{eq:cheb-Ck-case2}
\sum_{(i, j)\in S_t}\Ex[\bar{\chi}_{i}\cdot \bar{\chi}_{j}] 
\leq  {\ell \choose t} \cdot s^{2\ell -t} \cdot |C|\cdot   \left(\frac{1}{x}\right)^{2\ell -t} \leq \ell^t \cdot |C| \cdot \left(\frac{s}{x}\right)^{2\ell -t} \leq  \ell^\ell \cdot |C| \cdot   \left(\frac{s}{x}\right)^\ell =  \left( \frac{|C|^{1/\ell} \cdot \ell \cdot s}{x}\right)^\ell \leq (16\ell^2)^\ell\;,
\end{equation}

Thus,

\begin{equation}
\frac{\sum_{i,j} \Ex[\bar{\chi}_{i}\cdot \bar{\chi}_{j}]}{\theta^2\cdot \mu^2}
\leq  1/(\theta \mu) + (\ell (16\ell^2)^\ell)/(\theta \mu)^2 \leq 
1/ (8\ell)^\ell + (\ell (16\ell^2)^\ell)/(8\ell)^{2\ell} \leq 1/8 + \ell (1/2)^{2\ell} \leq 3/8\;,
\end{equation}
which concludes the proof of the lemma.
\end{proof}

%\begin{equation}\label{eq:cheb-Ck-caseX}
%\Ex[\chi_{i} \chi_{j}] =  \Ex\left[\sum_{h\in [\eta]} \psi^h_{i} \sum_{r \in [\eta]} \psi^r_{j} \right]\;,
%\end{equation}

\begin{theorem}\label{thm:Gen-F-ub}
There exists a one-sided error tester for $F$-freeness whose query complexity is $$O\left(|V(F)|^{2+1/\ell(F)} \cdot m \cdot \left(\frac{\alpha}{m }\right)^{1/\ell(F)} \cdot \left(\frac{1}{\epsilon}\right)^{1+2/\ell(F)}\right)$$ 
\end{theorem}

\begin{proof}
The completeness of the tester follows by construction. 
We next prove its soundness. Assume $G$ is $\epsilon$-far from being $F$-free.
%By Claim~\ref{lem:alpha}, w.p. at least $5/6$,
By Claim~\ref{clm:arb} and the setting of $t$, $|E(H(G, t))| \geq (1-(\epsilon/2)) |E(G)|$. 
Thus, by Claim~\ref{clm:disjoint}, there exist at least $\epsilon m/ (2t |F|)$ light-vertex-disjoint copies of $F$ in $H(G, t)$. \commentt{$|F|$ is $|V(F)|$, right?}
From each one of these copies we pick a subset of light vertices which covers the edges of this copy (since the copy is in $H(G, t)$ there exists a vertex cover which is composed of only light vertices).
Moreover, by the definition of $\ell(F)$, it suffices to pick $\ell(F)$ light vertices for each copy.
By ordering these vertices (arbitrarily) we obtain an $\ell(F)$-tuple of light vertices to each one of these copies.  
Let $C$ to be the set of $\ell(F)$-tuples corresponding to each one of the copies. According to Claim~\ref{clm:ver-num}, if we set $s$ to be at least $16 \ell(F) n / ((\epsilon m/ (2t |F|))^{1/\ell(F)})$ vertices, then w.h.c.p. there exists a copy of $F$ in $G$ such that we hit all its light vertices.
By Claim~\ref{clm:average}, w.h.p., the query complexity of the algorithm is bounded by $s \cdot 2\bar{d} = O\left(|V(F)|^{2+1/\ell(F)} \cdot m \cdot \left(\frac{\alpha}{m }\right)^{1/\ell(F)} \cdot \left(\frac{1}{\epsilon}\right)^{1+2/\ell(F)}\right)$, as desired.
\end{proof}

\begin{corollary}\label{cor:Ck-ub}
There exists a one-sided error tester for $C_k$-freeness whose query complexity is 
$$O\left(k^{2+(2/k))} \cdot m \cdot \left(\frac{\alpha}{m }\right)^{2/k} \cdot \left(\frac{1}{\epsilon}\right)^{1+(4/k)}\right)$$    
for even $k$, and 
$$O\left(k^{2+2/(k+1)} \cdot m \cdot \left(\frac{\alpha}{m }\right)^{2/(k+1)} \cdot \left(\frac{1}{\epsilon}\right)^{1+4/(k+1))}\right)$$    
for odd $k$.
\end{corollary}

The following bound is an improved bound for testing $C_k$-freeness for odd $k$ whenever $m = \Omega(\alpha^{(k+3)/2})$ and in particular when $\alpha$ is a constant.
\begin{theorem}\label{thm:Ck-ub-odd-k}
There exists a one-sided error tester for $C_k$-freeness for odd $k$ whose query complexity is 
$$O\left(k \cdot m \cdot (\alpha^2/m)^{\frac{2}{k-1}} \cdot \left( \frac{1}{\epsilon}\right)^{\frac{6}{k-1}} \right)$$
\end{theorem}

\begin{proof}
The completeness of the tester follows by construction. \commentt{need to say which tester exactly this is}
We next prove its soundness. Assume $G$ is $\epsilon$-far from being $F$-free.
By Claim~\ref{clm:arb} and the setting of $t$, $|E(H(G, t))| \geq (1-(\epsilon/2)) |E(G)|$. 
Thus, by Claim~\ref{clm:disjoint}, there exist at least $\epsilon m/ (2t k)$ light-vertex-disjoint copies of $C_k$ in $H(G, t)$.
From each one of these copies we pick a subset of light vertices which covers the edges of this copy (since the copy is in $H(G, t)$ there exists a vertex cover which is composed of only light vertices).
By the structure of $C_k$, it suffices to pick $\lceil k/2 \rceil$ light vertices from each copy.
Moreover, exactly one pair of these vertices are on the same edge of the corresponding copy $C_k$.
By excluding this pair and ordering the rest of these vertices (arbitrarily) , we obtain an $(k-3)/2)$-tuple of light vertices to each one of these copies.  
Let $X$ denote the set of these tuples.
We also consider the set of pairs of vertices we excluded from each copy.
Each such pair is an edge of a $C_k$. 
Denote this set of edges by $E$.
This set is of size $|X|$ which is at least $\epsilon m/ (2t k)$.
By Claim~\ref{clm:sampleedges}, w.h.p., we sample at least $c \epsilon s_1/ (2t k) $ edges from $E$ (where $c$ is a constant).
Let $C \subseteq X$ be the set of tuples corresponding to copies for which we sampled an edge from $E$.
By the above, w.h.p., $|C| \geq c \epsilon s_1/ (2t k) $. 
Conditioning on this event, according to Claim~\ref{clm:ver-num}, if $s_2$ is at least $16 k n / (c \epsilon s_1/ (2t k))^{2/(k-3)}$, then w.h.c.p. the algorithm will find a copy of $C_k$. 
By Claim~\ref{clm:average} and Claim~\ref{clm:sample}, the expected query complexity of the algorithm is bounded by $O(s_1 \cdot t + s_2 \cdot \bar{d}) = O\left(k \cdot m^{1- \frac{2}{k-1}} \cdot \alpha^{\frac{4}{k-1}} \cdot \left( \frac{1}{\epsilon}\right)^{\frac{6}{k-1}} \right)$, as desired.
\end{proof}

\begin{theorem}\label{thm:Gen-F-ub2}
There exists a one-sided error tester for $F$-freeness whose expected query complexity is $$O\left(|E(F)|^{1+1/|E(F)|} \cdot n\alpha \cdot \left(\frac{1}{m }\right)^{1/|E(F)|} \cdot \left(\frac{1}{\epsilon}\right)^{1+1/|E(F)|}\right)\;.$$ 
\end{theorem}

\begin{proof}
The completeness of the tester follows by construction. 
We next prove its soundness. Assume $G$ is $\epsilon$-far from being $F$-free.
Thus, there exist at least $\epsilon m/ (2|E(F)|)$ edge-disjoint copies of $F$ in $H(G, t)$.
According to Claim~\ref{clm:ver-num}, if we sample at least $16 |E(F)| m / ((\epsilon m/ (2|E(F)|))^{1/|E(F)|})$ edges, then w.h.c.p. there exists a copy of $F$ in $G$ such that we hit all its edges.
The expected query complexity of the algorithm is bounded by $h \cdot (t/\bar{d})$. By the setting of $t$ we obtain the desired result. 
\end{proof}

\begin{claim}\label{clm:average}
Let $S=\{v_1, \ldots, v_s\}$ be a uniform sample of size at least $s=(t/\bar{d})\cdot 12\ln(1/\delta)$ vertices (with repetitions), where $t$ is any threshold such that $|E(H(G, t))| \geq (1- \epsilon) m$.
Let $\chi_i=d(v_i)$ if $d(v_i)\leq t$ and $\chi_i=0$ otherwise.
Then with probability at least $1-\delta$, $\frac{1}{s}\sum_{i=1}^s\chi_i\leq 2\bar{d}$.
\end{claim}

\begin{proof}
It holds that $\Ex[\chi_i] \leq \frac{1}{n} \sum_{v \in V} d(v) = \bar{d}$.
%Also, let $\mathcal{L}$ denote the set of vertices with degree at most $\alpha/\eps$. In every graph with arboricity at most $\alpha$, there are at most $\eps m$ in the subgraph induced by $V \setminu \mathcal{L}$, and every edge that does not belong to that subgraph has at least one endpoint in $\mathcal{L}$. 
For any $\epsilon < \frac{1}{2}$, 
$\Ex[\chi_i]\geq
\frac{1}{n}\sum_{v \in L_t(G)} d(v) 
\geq \frac{1}{n}\cdot 2(1-\epsilon)m
\geq \bar{d}$.

Also, for every $i\in [s]$, $\chi_i\in [0,t]$ for every $i$. 
Let $x_i=\chi/t$ and $X=\sum_{i=1}^s x_i$.
Then the $x_i$'s are bounded independent   random variables in $[0,1]$ with $\Ex[X]\leq s\cdot \bar{d}/t$. Therefore, by the multiplicative Chernoff's bound
\[
\Pr\left[X > 2s\cdot \frac{\bar{d}}{t}\right] \leq  \Pr[X > 2s\cdot \Ex[X]] < \exp\left( -\frac{\Ex[X]}{3}\right) < \exp\left( -\frac{s\cdot (\bar{d}/4)/(\alpha/\eps)}{3}\right)< \delta.
\]
\end{proof}